% !TEX root = ../thesis.tex

\chapter{Knižnica libkarel}\label{ch:libkarel}

Existuje niekoľko variánt \emph{Robota Karla}, ktoré sa od seba líšia množstvom primitív a dostupných senzorov. Preto aj v rámci tejto práce sme sa rozhodli implementovať tri varianty knižnice:

\begin{enumerate}
    \item knižnicu \texttt{karel.h}, ktorá obsahuje len základné primitívy a senzory pre \emph{Robota Karla},
    \item knižnicu \texttt{superkarel.h}, ktorá obsahuje mnoho rozširujúcich príkazov a senzorov pre \emph{Robota Karla}, a
    \item knižnicu \texttt{bit.h}, ktorá je podľa \cite{Gregg2021} odvodená z knižnice \emph{Robota Karla} a vo výučbe ju na \emph{Stanfordskej univerzite} používal v úvodných kurzoch programovania \emph{CS106A}\footnote{\url{https://cs106a.stanford.edu/}} lektor \emph{Nick Parlante}.
\end{enumerate}

V tejto kapitole budú postupne predstavené hlavné stavebné bloky pre vytvorenie knižnice, formát súboru pre reprezentáciu sveta \emph{Robota Karla}, ako aj jednotlivé uvedené varianty.

globalne premenne

\section{Formát vstupného súboru}


\section{Knižnica karel.h}

ako vytvorit \texttt{loop()} a funkciu \texttt{turn\_on()} bez parametrov

Kompletná dokumentácia knižnice sa nachádza v prílohe \nameref{app:karel}.


\section{Knižnica superkarel.h}

Dokumentácia knižnice sa nachádza v prílohe \nameref{app:karel}.


\section{Knižnica bit.h}

Dokumentácia knižnice sa nachádza v prílohe \nameref{app:bit}.


\lstinputlisting[caption={Riešenie problému Schody},language=C]{listings/stairs.c}
